\section{The Fourier Transform}
\subsection{Definition and Simple Properties}

\begin{definition}
    For a integrable function \(\functiondomain{f}{\RR}{\CC}\), we define the \emph{Fourier transform} of \(f\) as the function \(\functiondomain{\hat{f}}{\RR}{\CC}\) given by
    \begin{equation}
        \hat{f}\pqty{\lambda} = \int_{\minusinfty}^{\plusinfty} f\pqty{x} e^{-i \lambda x} \dd{x},
    \end{equation}
    where \(\lambda \in \RR\).
\end{definition}

We often denote \(\fourier\) as the map mapping \(\functionmap{\fourier}{f}{\hat{f}}\), or more explicitly, we write the variables we take the Fourier transform with respect to, as in \(\fourier[x][\lambda] \bqty{f\pqty{x}} = \hat{f}\pqty{\lambda}\).

We assume that we work on sufficiently nice functions (e.g. on a Schwartz space) so that the Fourier transform is well-defined, and the relevant integrals converge. We also expect the \autoref{thm:riemann-lebesgue-lemma} (Riemann--Lebesgue lemma) to hold, i.e., the Fourier transform of an integrable function to decay for large \(\lambda\). For simplicity, we denote
\begin{equation}
    \int = \int_{\minusinfty}^{\plusinfty}
\end{equation}
in this section.

\begin{proposition}[Properties of the Fourier Transform]
    We have the following properties of the Fourier transform.

    \begin{itemize}
        \item \textbf{Differentiation and Multiplying by Polynomials.} We have
        \begin{subequations}
            \begin{align}
                \fourier[x][\lambda] \bqty{\pqty{\dv{x}}^k f\pqty{x}} &= \pqty{i \lambda}^k \hat{f}\pqty{\lambda}, \label{eq:fourier-diff}\\
                \fourier[x][\lambda] \bqty{x^k f\pqty{x}} &= i^k \pqty{\dv{\lambda}}^k \hat{f}\pqty{\lambda}. \label{eq:fourier-poly}
            \end{align}
        \end{subequations}
    
        \item \textbf{Shift and Phase Change.} We have
        \begin{subequations}
            \begin{align}
                \fourier[x][\lambda] \bqty{f\pqty{x - a}} &= e^{-i \lambda a} \hat{f}\pqty{\lambda}, \\
                \fourier[x][\lambda] \bqty{e^{- i a x} f\pqty{x}} &= \hat{f}\pqty{\lambda + a}.
            \end{align}
        \end{subequations}
    \end{itemize}
\end{proposition}

\begin{proof}
    We show the first property on differentiation and polynomials. For \eqref{eq:fourier-diff}, using repeated integration by parts, we have
    \begin{subequations}
        \begin{align}
            \fourier[x][\lambda] \bqty{\pqty{\dv{x}}^k f\pqty{x}} &= \int \dv[k]{x} f\pqty{x} e^{-i \lambda x} \dd{x} \\
            &= \int f\pqty{x} \pqty{- \dv{x}}^k e^{-i \lambda x} \dd{x} \\
            &= \int f\pqty{x} \pqty{i \lambda}^k e^{-i \lambda x} \dd{x} \\
            &= \pqty{i \lambda}^k \hat{f}\pqty{\lambda}.
        \end{align}
    \end{subequations}

    For \eqref{eq:fourier-poly}, we have
    \begin{subequations}
        \begin{align}
            i^k \pqty{\dv{\lambda}}^k \hat{f}\pqty{\lambda} &= i^k \pqty{\dv{\lambda}}^k \int f\pqty{x} e^{-i \lambda x} \dd{x} \\
            &= \pqty{i \dv{\lambda}}^k \int f\pqty{x} e^{-i \lambda x} \dd{x} \\
            &= \int f\pqty{x} \pqty{i \dv{\lambda}}^k e^{-i \lambda x} \dd{x} \\
            &= \int f\pqty{x} x^k e^{-i \lambda x} \dd{x} \\
            &= \fourier[x][\lambda] \bqty{x^k f\pqty{x}},
        \end{align}
    \end{subequations}
    by differentiating under the integral.
\end{proof}

In particular, note that \eqref{eq:fourier-diff} is extremely useful. Suppose we have a differential equation over \(x \in \RR\) of the form
\begin{equation}
    P\pqty{\dv{x}} y = f,
\end{equation}
where \(P\) is a polynomial in \(\dv*{x}\). Then, taking the Fourier transform, we have
\begin{equation}
    P\pqty{i \lambda} \hat{y}\pqty{\lambda} = \hat{f}\pqty{\lambda},
\end{equation}
which simply solves to
\begin{equation}
    \hat{y} \pqty{\lambda} = \frac{\hat{f}\pqty{\lambda}}{P\pqty{i \lambda}}.
\end{equation}

To find \(y = y\pqty{x}\), we would need to invert the Fourier transform.

\begin{proposition}[Inversion Formula for the Fourier Transform]
    We have the inversion formula for the Fourier transform, given by
    \begin{equation}
        f\pqty{x} = \frac{1}{2 \pi} \int \hat{f}\pqty{\lambda} e^{i \lambda x} \dd{\lambda}.
    \end{equation}
\end{proposition}

\begin{proof}
    We have
    \begin{subequations}
        \begin{align}
            \frac{1}{2\pi} \int_{-n}^{n} e^{i \lambda x} \hat{f}\pqty{\lambda} \dd{\lambda} &= \int f\pqty{y} \bqty{\frac{1}{2\pi} \int_{-n}^{n} e^{i \lambda \pqty{x - y}} \dd{\lambda}} \dd{y} \\
            &= \int f\pqty{y} \frac{\sin \pqty{n \pqty{x - y}}}{\pi \pqty{x - y}} \dd{y} \\
            &= \int f\pqty{x + y} \frac{\sin \pqty{n y}}{\pi y} \dd{y}.
        \end{align}
    \end{subequations}

    Recall from \textit{IA Differential Equations} that we have
    \begin{equation}
        \int \frac{\sin \pqty{n y}}{\pi y} \dd{y} = \int \frac{\sin y}{\pi y} \dd{y} = 1.
    \end{equation}

    Using this, we have
    \begin{subequations}
        \begin{align}
            \frac{1}{2\pi} \int_{-n}^{n} e^{i \lambda x} \hat{f}\pqty{\lambda} \dd{\lambda} - f\pqty{x} &= \int \frac{\sin \pqty{n y}}{\pi y} \bqty{f\pqty{x + y} - f\pqty{x}} \dd{y} \\
            &= \int \sin\pqty{ny} F\pqty{x, y} \dd{y},
        \end{align}
    \end{subequations}
    where we define the new function as
    \begin{equation}
        F\pqty{x, y} = \frac{f\pqty{x + y} - f\pqty{x}}{\pi y}.
    \end{equation}

    Using integration by parts, we may concludes that
    \begin{subequations}
        \begin{align}
            \int \sin\pqty{ny} F\pqty{x, y} \dd{y} &= \frac{1}{n} \int \cos\pqty{ny} \dv{y} F\pqty{x, y} \dd{y} \\
            &\to 0
        \end{align}
    \end{subequations}
    as \(n \to \infty\), by the \autoref{thm:riemann-lebesgue-lemma} (Riemann--Lebesgue Lemma).
\end{proof}

If we remove the limiting process, we have the \emph{formal} formula
\begin{equation}
    f\pqty{x} = \int f\pqty{y} \pqty{\frac{1}{2\pi} \int e^{i \lambda \pqty{x - y}} \dd{\lambda}} \dd{y},
\end{equation}
and the Dirac \(\diracdelta\) function therefore satisfies
\begin{equation}
    \label{eq:dirac-delta-fourier}%
    \diracdelta\pqty{x - y} = \frac{1}{2\pi} \int e^{i \lambda \pqty{x - y}} \dd{\lambda}.
\end{equation}

\begin{proposition}[Parseval's Theorem]
    We have
    \begin{equation}
        \int f\pqty{x} \conjugate{g\pqty{x}} \dd{x} = \frac{1}{2\pi} \int \hat{f}\pqty{\lambda} \conjugate{\hat{g}\pqty{\lambda}} \dd{\lambda}.
    \end{equation}
\end{proposition}

\begin{proof}
    We have
    \begin{subequations}
        \begin{align}
            \frac{1}{2\pi} \int \hat{f}\pqty{\lambda} \conjugate{\hat{g}\pqty{\lambda}} \dd{\lambda} &= \frac{1}{2\pi} \int \bqty{\int e^{-i \lambda y} f\pqty{y} \dd{y}} \bqty{\int e^{i \lambda z} \conjugate{g\pqty{x}} \dd{z}} \dd{\lambda} \\
            &= \iint f\pqty{y} \conjugate{g\pqty{x}} \bqty{\frac{1}{2\pi} \int e^{i \lambda \pqty{x - y}} \dd{\lambda}} \dd{\pqty{x, y}} \\
            &= \iint f\pqty{y} \conjugate{g\pqty{x}} \diracdelta\pqty{x - y} \dd{\pqty{x, y}} \\
            &= \int f\pqty{x} \conjugate{g\pqty{x}} \dd{x}.
        \end{align}
    \end{subequations}
\end{proof}

\begin{definition}[Convolution of Functions]
    For two functions \(\functiondomain{f, g}{\RR}{\CC}\), we define the \emph{convolution} of \(f\) and \(g\) as the function \(\functiondomain{f \convolution g}{\RR}{\CC}\) given by
    \begin{equation}
        f \convolution g\pqty{x} = \int f\pqty{y} g\pqty{x - y} \dd{y}.
    \end{equation}
\end{definition}

Note that \(f \convolution g = g \convolution f\).

\begin{proposition}[Convolution Theorem]
    \label{prop:convolution-theorem}%
    We have
    \begin{equation}
        \fourier[x][\lambda] \bqty{f \convolution g\pqty{x}} = \hat{f}\pqty{\lambda} \hat{g}\pqty{\lambda}.
    \end{equation}
\end{proposition}

\begin{proof}
    We have
    \begin{subequations}
        \begin{align}
            \fourier[x][\lambda] \bqty{f \convolution g\pqty{x}} &= \int e^{- i \lambda x} \bqty{\int f\pqty{x - y} g\pqty{y} \dd{y}} \dd{x} \\
            &= \int \bqty{\int f\pqty{x - y} g\pqty{y} e^{- i \lambda x} \dd{x}} \dd{y} \\
            &= \int \pqty{\int f\pqty{x} g\pqty{y} e^{- i \lambda \pqty{x + y}} \dd{x}} \dd{y} \\
            &= \int f\pqty{x} e^{- i \lambda x} \dd{x} \int g\pqty{y} e^{- i \lambda y} \dd{y} \\
            &= \hat{f}\pqty{\lambda} \hat{g}\pqty{\lambda}.
        \end{align}
    \end{subequations}
\end{proof}

\subsection{Important Examples of Fourier Transforms}

\paragraph{Exponentials.} Take \(f\pqty{x} = H\pqty{x} \exp\pqty{- \sigma x}\) where \(\Re\pqty{\sigma} > 0\), where \(H\pqty{x}\) is the Heaviside step function, then
\begin{equation}
    \hat{f}\pqty{\lambda} = \int_{0}^{\plusinfty} e^{-\pqty{\sigma + i \lambda}} x \dd{x} = \frac{1}{\sigma + i \lambda}.
\end{equation}

Note that \(\hat{f}\pqty{\lambda} \to 0\) as \(\abs{\lambda} \to \infty\), as expected from the Riemann--Lebesgue lemma. From the inversion formula, we may conclude that
\begin{subequations}
    \begin{align}
        H\pqty{x} e^{-\sigma x} &= \frac{1}{2\pi} \int \frac{e^{i \lambda x}}{\sigma + i \lambda} \dd{\lambda} \\
        &= \frac{1}{2 \pi i} \int \frac{e^{i \lambda x}}{\lambda - i \sigma} \dd{\lambda}.
    \end{align}
\end{subequations}

As we will see in \textit{IB Complex Analysis} or \textit{IB Complex Methods}, we can integrate the right-hand side using contour integration. We could also differentiate both sides by \(\pqty{-\pdv*{\sigma}}^k\) to obvain
\begin{equation}
    H\pqty{x} x^k e^{- \sigma x} = \frac{k!}{2 \pi} \int \frac{e^{i \pi x}}{\pqty{\sigma + i \lambda}^{k + 1}} \dd{\lambda}.
\end{equation}

\paragraph{Gaussians.} Let \(f\pqty{x} = \flatfrac{e^{- \flatfrac{x^2}{2}}}{\sqrt{2 \pi}}\). Directly computing the Fourier transform, we get
\begin{equation}
    \hat{f}\pqty{\lambda} = \int \frac{1}{\sqrt{2 \pi}} e^{- i \lambda x - \flatfrac{x^2}{2}} \dd{x}
\end{equation}
but this is hard to compute without knowing complex integration.

\small

By completing the square, we have
\begin{subequations}
    \begin{align}
        \hat{f}\pqty{\lambda} &= \frac{1}{\sqrt{2 \pi}} \int e^{- \flatfrac{x^2}{2} - i \lambda x} \dd{x} \\
        &= \frac{1}{\sqrt{2 \pi}} \int e^{- \flatfrac{\pqty{x + i \lambda}^2}{2}} e^{- \flatfrac{\lambda^2}{2}} \dd{x} \\
        &= e^{- \flatfrac{\lambda^2}{2}} \frac{1}{\sqrt{2 \pi}} \int e^{- \flatfrac{\pqty{x + i \lambda}^2}{2}} \dd{x} \\
        &= e^{- \flatfrac{\lambda^2}{2}},
    \end{align}
\end{subequations}
using the integration formula
\begin{equation}
    \int_{\minusinfty}^{\plusinfty} \exp\pqty{- \frac{\alpha}{2} \pqty{k - \gamma}^2} \dd{k} = \sqrt{\frac{2 \pi}{\alpha}},
\end{equation}
which holds for any \(\gamma\) and complex \(\alpha\) with \(\Re\pqty{\alpha} > 0\).

\normalsize

Instead, note that \(f\) satisfies the differential equation
\begin{equation}
    \pqty{\dv{x}} f = - x f,
\end{equation}
and by taking Fourier transforms of both sides, we have
\begin{equation}
    \pqty{i \lambda} \hat{f} = - \pqty{i \dv{\lambda}} \hat{f},
\end{equation}
which simplifies to
\begin{equation}
    \pqty{\dv{\lambda}} \hat{f} = - \lambda \hat{f}.
\end{equation}

This solves to \(\hat{f}\pqty{\lambda} = \hat{f} \pqty{0} e^{-\flatfrac{\lambda^2}{2}}\). But note \(\hat{f}\pqty{0} = \int f\pqty{x} \dd{x} = 1\), so \(\hat{f}\pqty{\lambda} = e^{-\flatfrac{\lambda^2}{2}}\). We also discovered that \(f\) is an eigenfunction of the Fourier transform, with eigenvalue \(\sqrt{2 \pi}\).

\paragraph{Dirac \(\diracdelta\) Function.} Consider the Dirac \(\diracdelta\) function, and note
\begin{equation}
    \hat{\delta} = \int e^{- i \lambda x} \diracdelta\pqty{x} \dd{x} = 1.
\end{equation}

USing the inversion formula, we get
\begin{equation}
    \diracdelta\pqty{x} = \frac{1}{2 \pi} \int e^{i \lambda x} \dd{\lambda},
\end{equation}
which is consistent with what we see in \eqref{eq:dirac-delta-fourier}.

\subsection{Vanishing Initial Value Problems}

Consider the initial value problem on \(t > 0\) with
\begin{subequations}
    \begin{align}
        Ly = f &\qfor t > 0, \\
        y = 0 &\qfor t = 0, \\
        \dot{y} = 0 &\qfor t = 0,
    \end{align}
\end{subequations}
where \(L\) is a second-order differential operator of the form
\begin{equation}
    L = \dv[2]{t} + \beta \dv{t} + \gamma.
\end{equation}

We would like to use Fourier transform, and hence we extend functions \(y, f\) to the region \(t < 0\) by setting them to be identically zero.

When doing Fourier transform with time variable, it is common to conduct \(\fourier[t][\omega]\). Then, we have
\begin{equation}
    \bqty{\pqty{i \omega}^2 + \beta \pqty{i \omega} + \gamma} \hat{y}\pqty{\omega} = \hat{f}\pqty{\omega},
\end{equation}
for all \(\omega \in \RR\). Hence,
\begin{equation}
    \hat{y} \pqty{\omega} = \frac{\hat{f}\pqty{\omega}}{P\pqty{\omega}} \qfor P\pqty{\omega} = \pqty{i \omega}^2 + \beta \pqty{i \omega} + \gamma.
\end{equation}

We can factorise \(P\pqty{\omega} = \pqty{i \omega + \sigma_1} \pqty{i \omega + \sigma_2}\) for some \(\sigma_1, \sigma_2\). We assume that \(\sigma_1 \neq \sigma_2\) and \(\Re\pqty{\sigma_{1, 2}} > 0\). Then, conducing partial fractions gives
\begin{subequations}
    \begin{align}
        \frac{1}{P\pqty{\omega}} &= \frac{1}{\pqty{i \omega + \sigma_1} \pqty{i \omega + \sigma_2}} \\ 
        &= \frac{1}{\sigma_2 - \sigma_1} \bqty{\frac{1}{i \omega + \sigma_1} - \frac{1}{i \omega + \sigma_2}} \\
        &= \frac{1}{\sigma_2 - \sigma_1} \fourier[t][\omega] \bqty{H\pqty{t} e^{- \sigma_1 t} - H\pqty{t} e^{- \sigma_2 t}},
    \end{align}
\end{subequations}
where we did indeed use the fact that \(\Re\pqty{\sigma_{1, 2}} > 0\)

The function defined as
\begin{equation}
    R\pqty{t} = H\pqty{t} \bqty{\frac{e^{- \sigma_1 t} - e^{- \sigma_2 t}}{\sigma_2 - \sigma_1}}
\end{equation}
is known as the \term{response function} associated with the differential operator \(L\). In the degenerate case where \(\sigma_1 = \sigma_2 = \sigma\), using detuning by setting \(\sigma_1 = \sigma\) and \(\sigma_2 = \sigma + \epsilon\) and taking the limit as \(\epsilon \to 0\), we have
\begin{equation}
    R\pqty{t} = H\pqty{t} t e^{- \sigma t},
\end{equation}
and we have
\begin{equation}
    \hat{y}\pqty{\omega} = \hat{R}\pqty{\omega} \hat{f}\pqty{\omega}.
\end{equation}

By \autoref{prop:convolution-theorem} (Convolution Theorem), we have
\begin{equation}
    y\pqty{t} = \pqty{R \convolution f} \pqty{t} = \int_{0}^{t} R\pqty{t - \tau} f\pqty{\tau} \dd{\tau},
\end{equation}
and therefore \(R\pqty{t - \tau} = G\pqty{t, \tau}\) is the Green's function associated with the IVP for \(L\). If we take an \term{impulsive forcing} \(f\pqty{t} = \diracdelta\pqty{t}\), then the solution is
\begin{equation}
    y\pqty{t} = \int_{0}^{t} R\pqty{t - \tau} \diracdelta\pqty{\tau} \dd{\tau} = R\pqty{t},
\end{equation}
hence the name response function.

\begin{example}[Damped Harmonic Oscillator]
    Consider a damped, simple harmonic oscillator with a forcing \(f\pqty{t}\) for \(t > 0\), with initial value problem (in suitable units)
    \begin{subequations}
        \begin{align*}
            \ddot{y} + 2 \gamma \dot{y} + y = f & \qfor t > 0, \\
            y = 0 & \qfor t = 0, \\
            \dot{y} = 0 & \qfor t = 0,
        \end{align*}
    \end{subequations}
    where \(\gamma > 0\) is the damping term. Taking Fourier transforms, we have
    \begin{equation}
        \hat{y}\pqty{\omega} = \frac{\hat{f}\pqty{\omega}}{P\pqty{\omega}} \qwhere P\pqty{\omega} = \pqty{i \omega + \gamma - \sqrt{\gamma^2 - 1}}\pqty{i \omega + \gamma + \sqrt{\gamma^2 - 1}}.
    \end{equation}

    Denote
    \begin{equation}
        \sigma_{\pm} = \gamma \pm \sqrt{\gamma^2 - 1},
    \end{equation}
    then note \(\Re\pqty{\sigma_{\pm}} > 0\) since \(\gamma > 0\). We conduct casework based on the value of \(\gamma\).

    \begin{subequations}
    \begin{enumerate}[label=(\roman*)]
        \item If \(\gamma > 1\), then
        \begin{equation}
            R\pqty{t} = H\pqty{t} \exp\pqty{- \gamma t} \frac{\sinh\pqty{t \sqrt{\gamma^2 - 1}}}{\sqrt{\gamma^2 - 1}}.
        \end{equation}

        This is known as \term{overdamping}.

        \item If \(0 < \gamma < 1\), then
        \begin{equation}
            R\pqty{t} = H\pqty{t} \exp\pqty{- \gamma t} \frac{\sin\pqty{t \sqrt{1 - \gamma^2}}}{\sqrt{1 - \gamma^2}}.
        \end{equation}

        This is known as \term{underdamping}.

        \item If \(\gamma = 1\), then
        \begin{equation}
            R\pqty{t} = H\pqty{t} t e^{- t}.
        \end{equation}

        This is known as \term{critical damping}.
    \end{enumerate}
    \end{subequations}
\end{example}